% GENERATED FILE. Edit canonical lessons, then run npm run generate:books.
% canonical-source-hash: fnv1a64:5744602f82480869
% canonical-lessons: HI-C461-sarminda, HI-C461-sakucana, HI-C461-hickicana, HI-C461-lajana, HI-C461-dhith

\chapter{Shyness and Shame}
\label{ch:hi-shyness-and-shame}
\hlchaptermodality{461}

\begin{chapteropening}
\textbf{By the end of this chapter:} \emph{I can say that someone felt shy or ashamed.}

Complete the last lesson of chapter 461: I can say that someone felt shy or ashamed.
\end{chapteropening}

\section[\texorpdfstring{\hi{शर्मिंदा} (śarmindā)}{शर्मिंदा (śarmindā)}]{\hi{शर्मिंदा} (śarmindā) — ashamed; embarrassed}

\label{lesson:HI-C461-sarminda}



\begin{quote}
Before the new one: say the Hindi for desolate, then the Hindi for a longing.
\end{quote}

\subsection*{You'll want to know: \hi{शर्मिंदा}}
\textbf{\hi{शर्मिंदा}} — \emph{śarmindā} — \textquotedblleft{}ashamed; embarrassed\textquotedblright{}.

\begin{quote}
\hi{अपनी} \hi{ग़लती} \hi{पर} \hi{अरुण} \hi{शर्मिंदा} \hi{था।}
\end{quote}

\begin{quote}
\emph{apnī ġaltī par Aruṇ śarmindā thā.}
\end{quote}

\begin{quote}
Arun was ashamed of his mistake.
\end{quote}

\begin{cousinweb}
\textbf{\hi{शर्मिंदा}} (\emph{śarmindā}) is usually traced to Persian \emph{sharminda}, \textquotedblleft{}ashamed\textquotedblright{}, the adjective beside the noun that gave Hindi \emph{śarm}.
\end{cousinweb}

\begin{grammarlens}[title={What you've built}]
One more word for describing what you lived and felt.
\end{grammarlens}

\subsection*{Your turn}
\begin{itemize}
\raggedright
  \item \emph{Say it:} \emph{śarmindā}
  \item \emph{Say it:} \emph{śarmindā}, once more
  \item \emph{Say it:} \emph{apnī ġaltī par Aruṇ śarmindā thā.} — Arun was ashamed of his mistake.
  \item \emph{Recall:} say the Hindi for desolate, then the Hindi for a longing, then say \emph{śarmindā} again
\end{itemize}

\begin{tcolorbox}[breakable,colback=teal!4,colframe=teal!35!black,title={Before you move on}]
What does \hi{शर्मिंदा} mean? (\textquotedblleft{}Ashamed; embarrassed\textquotedblright{}.) Say it once more.
\end{tcolorbox}
\section[\texorpdfstring{\hi{सकुचाना} (sakucānā)}{सकुचाना (sakucānā)}]{\hi{सकुचाना} (sakucānā) — to shrink back shyly}

\label{lesson:HI-C461-sakucana}



\begin{quote}
Before the new one: say the Hindi for a longing, then the Hindi for ashamed.
\end{quote}

\subsection*{You'll want to know: \hi{सकुचाना}}
\textbf{\hi{सकुचाना}} — \emph{sakucānā} — \textquotedblleft{}to shrink back shyly\textquotedblright{}.

\begin{quote}
\hi{पहली} \hi{मुलाक़ात} \hi{में} \hi{मीरा} \hi{थोड़ा} \hi{सकुचाई।}
\end{quote}

\begin{quote}
\emph{pahlī mulāqāt meṁ Mīrā thoṛā sakucāī.}
\end{quote}

\begin{quote}
At the first meeting Mira was a little shy and held back.
\end{quote}

Past (she …): \hi{वह} \hi{सकुचाई}, \emph{vah sakucāī}.

\begin{cousinweb}
\textbf{\hi{सकुचाना}} (\emph{sakucānā}) is inherited: CDIAL lists Hindi \emph{sakucnā}, \textquotedblleft{}to shrink, become narrow or tight\textquotedblright{}, under Sanskrit \emph{saṁkucyatē}, \textquotedblleft{}shrinks\textquotedblright{} (CDIAL 12824); \emph{sakucānā} is its common form for a person who shrinks back out of shyness.
\end{cousinweb}

\begin{grammarlens}[title={What you've built}]
One more word for describing what you lived and felt.
\end{grammarlens}

\subsection*{Your turn}
\begin{itemize}
\raggedright
  \item \emph{Say it:} \emph{sakucānā}
  \item \emph{Say it:} \emph{sakucānā}, once more
  \item \emph{Say it:} \emph{pahlī mulāqāt meṁ Mīrā thoṛā sakucāī.} — At the first meeting Mira was a little shy and held back.
  \item \emph{Recall:} say the Hindi for a longing, then the Hindi for ashamed, then say \emph{sakucānā} again
\end{itemize}

\begin{tcolorbox}[breakable,colback=teal!4,colframe=teal!35!black,title={Before you move on}]
What does \hi{सकुचाना} mean? (\textquotedblleft{}To shrink back shyly\textquotedblright{}.) Say it once more.
\end{tcolorbox}
\section[\texorpdfstring{\hi{हिचकिचाना} (hickicānā)}{हिचकिचाना (hickicānā)}]{\hi{हिचकिचाना} (hickicānā) — to hesitate}

\label{lesson:HI-C461-hickicana}



\begin{quote}
Before the new one: say the Hindi for ashamed, then the Hindi for to shrink back shyly.
\end{quote}

\subsection*{You'll want to know: \hi{हिचकिचाना}}
\textbf{\hi{हिचकिचाना}} — \emph{hickicānā} — \textquotedblleft{}to hesitate\textquotedblright{}.

\begin{quote}
\hi{सवाल} \hi{पूछने} \hi{से} \hi{पहले} \hi{अरुण} \hi{थोड़ा} \hi{हिचकिचाया।}
\end{quote}

\begin{quote}
\emph{savāl pūchne se pahle Aruṇ thoṛā hickicāyā.}
\end{quote}

\begin{quote}
Arun hesitated a little before asking his question.
\end{quote}

Past (he …): \hi{वह} \hi{हिचकिचाया}, \emph{vah hickicāyā}.

\begin{cousinweb}
\textbf{\hi{हिचकिचाना}} (\emph{hickicānā}) is probably built on \emph{hicaknā}, \textquotedblleft{}to draw back\textquotedblright{}, which CDIAL lists under \emph{*hicc}, a form it reconstructs (CDIAL 14083); the doubled \emph{-cki-} paints the stop-and-start of hesitation.
\end{cousinweb}

\begin{grammarlens}[title={What you've built}]
One more word for describing what you lived and felt.
\end{grammarlens}

\subsection*{Your turn}
\begin{itemize}
\raggedright
  \item \emph{Say it:} \emph{hickicānā}
  \item \emph{Say it:} \emph{hickicānā}, once more
  \item \emph{Say it:} \emph{savāl pūchne se pahle Aruṇ thoṛā hickicāyā.} — Arun hesitated a little before asking his question.
  \item \emph{Recall:} say the Hindi for ashamed, then the Hindi for to shrink back shyly, then say \emph{hickicānā} again
\end{itemize}

\begin{tcolorbox}[breakable,colback=teal!4,colframe=teal!35!black,title={Before you move on}]
What does \hi{हिचकिचाना} mean? (\textquotedblleft{}To hesitate\textquotedblright{}.) Say it once more.
\end{tcolorbox}
\section[\texorpdfstring{\hi{लजाना} (lajānā)}{लजाना (lajānā)}]{\hi{लजाना} (lajānā) — to blush; to be bashful}

\label{lesson:HI-C461-lajana}



\begin{quote}
Before the new one: say the Hindi for to shrink back shyly, then the Hindi for to hesitate.
\end{quote}

\subsection*{You'll want to know: \hi{लजाना}}
\textbf{\hi{लजाना}} — \emph{lajānā} — \textquotedblleft{}to blush; to be bashful\textquotedblright{}.

\begin{quote}
\hi{सब} \hi{लोगों} \hi{के} \hi{सामने} \hi{अपना} \hi{नाम} \hi{सुनकर} \hi{मीरा} \hi{लजा} \hi{गई।}
\end{quote}

\begin{quote}
\emph{sab logoṁ ke sāmne apnā nām sunkar Mīrā lajā gaī.}
\end{quote}

\begin{quote}
Mira blushed when she heard her name in front of everyone.
\end{quote}

Past (she …): \hi{वह} \hi{लजा} \hi{गई}, \emph{vah lajā gaī}.

\begin{cousinweb}
\textbf{\hi{लजाना}} (\emph{lajānā}) is a tadbhava: CDIAL lists Hindi \emph{lajānā}, \textquotedblleft{}to shame, be ashamed\textquotedblright{}, under Sanskrit \emph{lajjátē}, \textquotedblleft{}is ashamed\textquotedblright{} (CDIAL 10909).
\end{cousinweb}

\begin{grammarlens}[title={What you've built}]
One more word for describing what you lived and felt.
\end{grammarlens}

\subsection*{Your turn}
\begin{itemize}
\raggedright
  \item \emph{Say it:} \emph{lajānā}
  \item \emph{Say it:} \emph{lajānā}, once more
  \item \emph{Say it:} \emph{sab logoṁ ke sāmne apnā nām sunkar Mīrā lajā gaī.} — Mira blushed when she heard her name in front of everyone.
  \item \emph{Recall:} say the Hindi for to shrink back shyly, then the Hindi for to hesitate, then say \emph{lajānā} again
\end{itemize}

\begin{tcolorbox}[breakable,colback=teal!4,colframe=teal!35!black,title={Before you move on}]
What does \hi{लजाना} mean? (\textquotedblleft{}To blush; to be bashful\textquotedblright{}.) Say it once more.
\end{tcolorbox}
\section[\texorpdfstring{\hi{ढीठ} (ḍhīṭh)}{ढीठ (ḍhīṭh)}]{\hi{ढीठ} (ḍhīṭh) — brazen; cheeky}

\label{lesson:HI-C461-dhith}



\begin{quote}
Before the new one: say the Hindi for to hesitate, then the Hindi for to blush.
\end{quote}

\subsection*{You'll want to know: \hi{ढीठ}}
\textbf{\hi{ढीठ}} — \emph{ḍhīṭh} — \textquotedblleft{}brazen; cheeky\textquotedblright{}.

\begin{quote}
\hi{वह} \hi{बच्चा} \hi{बहुत} \hi{ढीठ} \hi{है}, \hi{डाँटने} \hi{पर} \hi{भी} \hi{हँसता} \hi{है।}
\end{quote}

\begin{quote}
\emph{vah baccā bahut ḍhīṭh hai, ḍā̃ṭne par bhī hãstā hai.}
\end{quote}

\begin{quote}
That child is very cheeky; he laughs even when he's scolded.
\end{quote}

\begin{cousinweb}
\textbf{\hi{ढीठ}} (\emph{ḍhīṭh}) is a tadbhava: CDIAL lists Hindi \emph{ḍhīṭh}, \textquotedblleft{}bold, impudent\textquotedblright{}, under Sanskrit \emph{dhṛṣṭá}, \textquotedblleft{}bold\textquotedblright{} (CDIAL 6875).
\end{cousinweb}

\begin{grammarlens}[title={What you've built}]
One more word for describing what you lived and felt.
\end{grammarlens}

\subsection*{Your turn}
\begin{itemize}
\raggedright
  \item \emph{Say it:} \emph{ḍhīṭh}
  \item \emph{Say it:} \emph{ḍhīṭh}, once more
  \item \emph{Say it:} \emph{vah baccā bahut ḍhīṭh hai, ḍā̃ṭne par bhī hãstā hai.} — That child is very cheeky; he laughs even when he's scolded.
  \item \emph{Recall:} say the Hindi for to hesitate, then the Hindi for to blush, then say \emph{ḍhīṭh} again
\end{itemize}

\begin{tcolorbox}[breakable,colback=teal!4,colframe=teal!35!black,title={Before you move on}]
What does \hi{ढीठ} mean? (\textquotedblleft{}Brazen; cheeky\textquotedblright{}.) Say it once more.
\end{tcolorbox}
